\documentclass[11pt]{article}
\usepackage[T1]{fontenc}
\usepackage{lmodern,amsmath,amssymb,amsthm,geometry}
\usepackage[hidelinks]{hyperref}
\geometry{margin=1in}
\newtheorem{proposition}{Proposition}
\newtheorem{lemma}{Lemma}
\newtheorem{corollary}{Corollary}
\title{A five-witness quadratic presentation of the semilocal anchor}
\author{}
\date{November 6, 2026}
\begin{document}
\maketitle
\noindent
This note proves a replacement of one private module in the released
12,504-variable quartic. The polynomial projection equivalence of the
replacement is exact. The conclusion that the full polynomial defines
$\mathbb Z$ retains the arithmetic hypotheses and review status of the
upstream manuscript.

The old module asks, for $i=0,1$, that $Ca_i=f_i$ and that
\begin{align*}
 b_{i0}^2-2b_{i1}^2-3b_{i2}^2+6b_{i3}^2&=1,\\
 (a_i-2b_{i0})^2-8c_{i1}^2-12c_{i2}^2+24c_{i3}^2&=4.
\end{align*}
The norm-one trace-sum theorem for the quaternion algebra $(2,3)$,
which is ramified at $2,3$ and split at infinity, identifies this
seven-witness predicate in $a_i$ with
$\mathbb Z_{(2)}\cap\mathbb Z_{(3)}$; see
\cite[Proposition 6]{Koenigsmann2016}.
The second equation is the norm-one equation for a quaternion whose
scalar part is $(a_i-2b_{i0})/2$, so the trace conventions match exactly.
The common $C$ is nonzero because the retained constant chain starts
from a nonzero integer and repeatedly squares it.
\section{An exact replacement for the semilocal anchor}

Write
\[
 R=\mathbb Z_{(2)}\cap\mathbb Z_{(3)}
   =\{t\in\mathbb Q:v_2(t)\ge0,\ v_3(t)\ge0\},
\]
with $v_p(0)=+\infty$. The following specialization of classical
quaternion methods is particularly economical. Daans proves three-quantifier
existential definability for every finite intersection of valuation rings
in a global field; he traces the corresponding rational cases back to
Robinson \cite[Proposition 4.2 and Remark 4.3]{Daans2024}.
Here all constants and all exceptional cases are made explicit.

\begin{proposition}\label{prop:explicit-semilocal}
For every $t\in\mathbb Q$,
\[
 t\in R\quad\Longleftrightarrow\quad
 \exists u,v,w\in\mathbb Q:\quad
 u^2+6v^2-2w^2-3t^2+2=0.
\]
\end{proposition}
\begin{proof}
Consider the quaternion algebra $B=(2,3)_{\mathbb Q}$ with generators
$i^2=2$, $j^2=3$, $ij=-ji$. Its ramification set consists exactly of
$2$ and $3$. Indeed its Hilbert symbol at $3$ equals the Legendre symbol
$(2/3)=-1$, its Hilbert symbol at $2$ is $-1$, it splits at every other
finite place because both parameters are units there, and it splits at
infinity because both parameters are positive. In particular $B$ is a
division algebra over $\mathbb Q$.

We use the standard quadratic embedding criterion for quaternion algebras:
a quadratic field $L/\mathbb Q$ embeds in $B$ if and only if every ramified
place of $B$ remains a field in $L$. This is the quadratic case of the
Albert--Brauer--Hasse--Noether theorem and the central simple algebra
embedding criterion; equivalently, extending scalars to $L$ kills the
local invariant $1/2$ precisely when the local extension has degree two.
See \cite[Chapter 14]{Voight2021}.

Suppose first that $t\ne0$, and put
\[
 \delta=1-\frac{2}{3t^2}.
\]
This number is nonzero: $t^2=2/3$ has no rational solution.
At $2$, put $m=v_2(t)$. If $m>0$, then $v_2(\delta)=1-2m$ is odd,
so $\delta$ is not a square in $\mathbb Q_2$. If $m=0$, then
$t^2\equiv1\pmod8$ and $\delta\equiv3\pmod8$, again a nonsquare.
If $m<0$, then $\delta\in1+8\mathbb Z_2$, hence it is a square in
$\mathbb Q_2$. At $3$, if $m=v_3(t)\ge0$ then
$v_3(\delta)=-1-2m$ is odd; if $m<0$, then
$\delta\in1+3\mathbb Z_3$ and Hensel's lemma makes it a square.
Consequently
\[
 \delta\notin\mathbb Q_2^{\times2},\quad
 \delta\notin\mathbb Q_3^{\times2}
 \quad\Longleftrightarrow\quad t\in R.
\]
The embedding criterion shows that this is equivalent to the existence
of a pure quaternion $xi+yj+zij$ with square $\delta$, that is,
\[
 2x^2+3y^2-6z^2=\delta.
\]
This equivalence also excludes the case in which $\delta$ is a rational
square: a division quaternion algebra cannot contain a noncentral element
whose square is a nonzero rational square, since
$(\alpha-s)(\alpha+s)=0$ would force $\alpha=\pm s$.
Multiplying the displayed equation by $3t^2$ and setting
$u=3ty$, $v=tx$, $w=3tz$ gives the equation in the statement.
These substitutions are invertible for $t\ne0$.
Finally $t=0$ belongs to $R$ and is represented by $(u,v,w)=(0,0,1)$.
\end{proof}

\begin{lemma}\label{lem:two-semilocal-inputs}
Let $F(A,B)=A^2+AB-B^2$. For every $a,b\in\mathbb Q$ and $p=2,3$,
\[
 v_p(F(a,b))=2\min\{v_p(a),v_p(b)\}.
\]
In particular $(a,b)\in R^2$ if and only if $F(a,b)\in R$.
\end{lemma}
\begin{proof}
The assertion is immediate for $a=b=0$. Otherwise scale $a,b$ by a
common power of $p$ until both are integral and at least one is a unit.
The reduction of $A^2+AB-B^2$ has no nontrivial zero over either
$\mathbb F_2$ or $\mathbb F_3$: over $\mathbb F_2$ one checks the
three nonzero pairs, and over $\mathbb F_3$ the discriminant $5$
is a nonsquare. Thus the scaled value is a unit, proving the equality.
\end{proof}

\begin{corollary}[Five-variable replacement for the anchor]
\label{cor:five-anchor}
Fix rational inputs $f_0,f_1,C$ with $C\ne0$. The condition
$f_0/C,f_1/C\in R$ is equivalent to the existence of five rational
variables $r,D,u,v,w$ satisfying the three quadratic equations
\begin{align*}
 r&=f_0^2+f_0f_1-f_1^2,\\
 D&=C^2,\\
 u^2+6v^2-2w^2-3r^2+2D^2&=0.
\end{align*}
\end{corollary}
\begin{proof}
By Lemma~\ref{lem:two-semilocal-inputs}, the desired condition is
$r/C^2\in R$, with $r=F(f_0,f_1)$. Apply
Proposition~\ref{prop:explicit-semilocal} to $t=r/C^2$ and multiply
its equation by $C^4$. Multiplication of each of its three witnesses
by $C^2$ is a bijection of $\mathbb Q$, giving the stated equation.
\end{proof}

The released anchor uses two scaling variables and two copies of a
seven-witness trace-sum condition, hence sixteen witnesses in this module.
Corollary~\ref{cor:five-anchor} replaces that module by five witnesses.
It therefore saves eleven variables before any other change and retains
quadratic residuals, so their sum of squares remains quartic. This count
concerns the specified module; an implementation should report the exact
final global count after removing any unused lifting definitions.
The three displayed residuals have integer coefficients of absolute value
at most six. This is an exact replacement, with no change to the arithmetic
meaning or strength of the anchor.

The same method handles arbitrary finite tuples. For a finite nonempty
set $S$ of rational primes and a positive integer $m$, choose a monic
polynomial $g\in\mathbb Z[T]$ of degree $m$ whose reductions are
irreducible modulo every $p\in S$; finite-field irreducible polynomials
and the Chinese remainder theorem supply one. If $\alpha$ is a root, let
\[
 N_m(X_1,\ldots,X_m)=
 N_{\mathbb Q(\alpha)/\mathbb Q}
 (X_1+X_2\alpha+\cdots+X_m\alpha^{m-1}).
\]
For $p\in S$ the extension is unramified and inert at $p$, and its
power basis is an integral basis locally. Therefore
\[
 v_p(N_m(\mathbf x))=m\min_i v_p(x_i).
\]
This proves, with no new existential variables, that all coordinates of
$\mathbf x$ belong to $\bigcap_{p\in S}\mathbb Z_{(p)}$ if and only if
$N_m(\mathbf x)$ does. It is the norm-form version of
\cite[Proposition 4.8]{Daans2024}. In the present two-prime, two-input
case the polynomial $T^2-T-1$ gives precisely the form $F$ above.


\begin{thebibliography}{9}
\bibitem{Daans2024} N. Daans, \emph{Universally defining $\mathbb Z$
in $\mathbb Q$ with $10$ quantifiers}, Journal of the London
Mathematical Society \textbf{109} (2024), no. 2, e12864.
\url{https://doi.org/10.1112/jlms.12864}.
\bibitem{Koenigsmann2016} J. Koenigsmann, \emph{Defining $\mathbb Z$
in $\mathbb Q$}, Annals of Mathematics \textbf{183} (2016), 73--93.
\url{https://doi.org/10.4007/annals.2016.183.1.2}.
\bibitem{Voight2021} J. Voight, \emph{Quaternion algebras},
Graduate Texts in Mathematics \textbf{288}, Springer, 2021.
\url{https://doi.org/10.1007/978-3-030-56694-4}.
\end{thebibliography}
\end{document}
